\cotefile{157-5}
\scanpage{261}\mspage{177}
% archive sheet number: 260
\[
\left\{\begin{aligned}
M_{\pi}&\subset\mathcal F^{\unsure{\pi|P}}_{\deleted{\unsure{\pi}}}
  &&(\text{si }\pi=\pi^{\pi_0}),\\
\text{donc }\ M_{\pi}&=\mathcal F_{\pi}
  &&\text{si \inserted{de plus} la \unsure{fcg|fag} est \uline{accessible}.}
\end{aligned}\right.
\tag*{(4.7.9.1)}
\]
\deleted{\illegible{1}}

\uline{Dém} \unsure{Appliquant|Appliquons} la condition c) des fcg à
un $F\in\operatorname{Ob}M_\pi$, on trouve que $\deleted{\unsure{F}}F$ est facteur
direct d'un $\deleted{\unsure{E}}F'\in\operatorname{Ob}\mathcal F_\pi$. Comme $\mathcal F^\pi_{\deleted{\unsure{\pi}}}$ est
stable par $\varinjlim$ \unsure{grdes|gdes} devant $\pi_0$ et de cardinal
$\leq\pi$, et que \uline{de plus} $\pi>\pi_0$, on conclut que

\uline{Lemme 4.7.9.1} : Si $\pi>\pi_0$,$^{*}$ $\mathcal F^\pi_{\deleted{\unsure{\pi}}}$ est stable par
facteurs directs,

% slanted
\margin{Mais $\mathcal F^{\pi_0}$ n'est pas forcément stable par facteurs directs, et comme on veut \unsure{retrouver} $C$ prenant $M=\operatorname{Ind}_{\pi_0}(C)$, $C$ une petite catégorie, on \unsure{ne voudra|voudrait} \unsure{pas|prendre} \unsure{ds} la filtration \unsure{cardinale|considérée}, \unsure{prendre} les $\mathrm{Filt}'^{\pi}$ --- ou, \unsure{notons} que $\mathrm{Filt}'^{\pi_0}=\mathrm{Filt}^{\pi_0}$ \unsure{S.} $\unsure{I\simeq C}$\deleted{\unsure{est tout petit}}, ou \unsure{quiconque|toute} cat. \deleted{\unsure{petite}} \unsure{grde} devant $\pi_0$ et de cardinal $\leq\pi_0$ est \unsure{cofinale} à $e$.}
il résulte que $F\in\mathcal F^\pi$, d'où (4.7.9.1).

\uline{Corollaire} \deleted{\unsure{4.}}4.7.10. Supposons la fcg accessible.
Pour tout cardinal $\pi\geq\pi_0$, on a
\[
\mathcal F^\pi_{\deleted{\unsure{\pi}}}\subset M_\pi\subset\mathcal F^{\pi^{\pi_0}}
\quad(=M^{\unsure{P|\pi}}_{\pi^{\pi_0}})
\tag*{(4.7.10.1)}
\]

En effet, $\mathcal F_\pi\subset M_\pi$ par définition de fcg
accessible, et $M_\pi\subset\deleted{\mathcal F^{\pi^{\pi_0}}}M_{\pi^{\pi_0}}$ (puisque
$\pi\leq\pi^{\pi_0}$), or $M_{\pi^{\pi_0}}=\mathcal F^{(\pi^{\pi_0})}$ en vertu de
(4.7.9.1) appliqué à $\pi'=\pi^{\pi_0}$ \unsure{ald} $\pi$ (vu que
$\pi'^{\pi_0}=\pi'$), d'où (4.7.10.1). On notera
que cette inclusion généralise (4.7.9.1),
elle donne $\mathcal F^\pi_{\deleted{\unsure{\pi}}}=M_\pi$ si $\pi=\pi^{\pi_0}$.

4.7.11. \deleted{\unsure{Il}} Il \inserted{me} semble que les seules \deleted{\unsure{condi}} filtrations,
cardinales \inserted{\unsure{gén.}} utiles sont celles qui sont \uline{accessibles}.
\deleted{\unsure{Et}} Dans la suite je vais me borner à ce
cas. Comme \inserted{alors} $\mathcal F_\pi\subset M_\pi$, on en conclut
